% =============================================================================
% Chapter 3: System Design and Methodology
% =============================================================================

\chapter{System Design and Methodology}
\label{ch:design}

The core challenge this thesis addresses is a retrieval problem with an unusual input: a natural language task description---a bug report, a feature request, an ambiguous question about behavior---must be mapped to the small subset of files in a large repository that are actually relevant. \Cref{ch:background} surveyed the available building blocks: sparse lexical search captures surface-level name matching; static dependency graphs encode structural relationships that text search cannot see; and Personalized PageRank (PPR) propagates query-specific relevance through graph structure. Each technique addresses a different dimension of the problem, and no single method handles all of them.

This chapter presents CodeGraph, a retrieval pipeline that assembles these building blocks into a coherent system. The guiding insight is a \emph{division of labor}: a multi-signal seed selection stage uses lexical methods to place the random walk in the right region of the graph, then PPR exploits structural coherence to rank results within that region. Lexical and structural evidence complement rather than compete with each other. Every design decision is described here in enough detail that a reader could reimplement the system from this chapter alone.

The chapter follows the data as it flows through the pipeline. \Cref{sec:architecture} introduces the five-stage architecture and its central design principle. \Cref{sec:parsing} examines how source code is parsed into typed entities. \Cref{sec:graph-construction} shows how those entities become a weighted property graph. \Cref{sec:seed-selection} addresses the translation of a natural language task into a set of PPR seed nodes. \Cref{sec:ppr-config} presents the PPR configuration and the rationale behind key parameter choices. \Cref{sec:post-processing} covers the post-processing steps that produce the final output. \Cref{sec:mcp-integration} explains how the pipeline is exposed to AI agents through the Model Context Protocol. Finally, \cref{sec:design-limitations} discusses the design-level limitations and trade-offs.


% =============================================================================
% 3.1 Architecture Overview
% =============================================================================
\section{Architecture Overview}
\label{sec:architecture}

CodeGraph is organized as a five-stage pipeline. Each stage consumes the output of the previous stage and produces a well-defined intermediate representation. The five stages are: \emph{source code parsing}, \emph{graph construction}, \emph{seed selection}, \emph{PPR-based ranking}, and \emph{context delivery}. \Cref{fig:pipeline-architecture} illustrates the full pipeline with its intermediate data types and the division between offline indexing and online retrieval.

\begin{figure}[H]
    \centering
    \begin{tikzpicture}[
            % Stage boxes with drop shadow effect
            stage/.style={draw, rounded corners=6pt, minimum width=3.4cm, minimum height=1.2cm,
                    text width=3.1cm, align=center, font=\small\bfseries,
                    line width=0.6pt},
            % Data labels between stages
            datalabel/.style={font=\scriptsize, text width=3.0cm, align=left,
                    draw=gray!30, rounded corners=2pt, fill=white,
                    inner sep=4pt},
            % Input/output boxes
            io/.style={draw, rounded corners=4pt, fill=gray!8, minimum width=3.4cm,
                    minimum height=0.9cm, font=\small, align=center, line width=0.5pt},
            % Arrow style
            arr/.style={-{Stealth[length=6pt, width=4pt]}, thick, color=gray!70},
            dataarr/.style={-{Stealth[length=4pt, width=3pt]}, thin, color=gray!40, dashed},
            node distance=0.7cm
        ]
        % === Left column: pipeline stages ===
        \node[io] (src) {Source Repository};

        \node[stage, fill=blue!12, draw=blue!40, below=of src] (parse)
        {\textcolor{blue!70}{Stage 1}\\Source Code\\Parsing};

        \node[stage, fill=green!10, draw=green!40, below=of parse] (graph)
        {\textcolor{green!50!black}{Stage 2}\\Graph\\Construction};

        \node[stage, fill=orange!12, draw=orange!50, below=of graph] (seed)
        {\textcolor{orange!70!black}{Stage 3}\\Seed\\Selection};

        \node[stage, fill=orange!18, draw=orange!60, below=of seed] (ppr)
        {\textcolor{orange!70!black}{Stage 4}\\PPR\\Ranking};

        \node[stage, fill=purple!10, draw=purple!40, below=of ppr] (deliver)
        {\textcolor{purple!60!black}{Stage 5}\\Context\\Delivery};

        \node[io, below=of deliver] (agent) {AI Agent};

        % === Right column: intermediate data ===
        \node[datalabel, right=1.0cm of parse.east, anchor=west] (d1)
        {\textcolor{blue!60}{$\blacktriangleright$} Entity Model\\
            \footnotesize File, Class, Function,\\Method entities};

        \node[datalabel, right=1.0cm of graph.east, anchor=west] (d2)
        {\textcolor{green!50!black}{$\blacktriangleright$} Property Graph\\
            \footnotesize Typed nodes +\\weighted edges in Neo4j};

        \node[datalabel, right=1.0cm of seed.east, anchor=west] (d3)
        {\textcolor{orange!70!black}{$\blacktriangleright$} Personalization Vector\\
            \footnotesize $\vec{s} \in \mathbb{R}^{|V|}$, $\|\vec{s}\|_1 = 1$};

        \node[datalabel, right=1.0cm of ppr.east, anchor=west] (d4)
        {\textcolor{orange!70!black}{$\blacktriangleright$} Ranked File List\\
            \footnotesize Top-$k$ files by\\PPR score};

        % === Task description input ===
        \node[io, left=2.0cm of seed] (task) {Task\\Description};

        % === Vertical arrows ===
        \draw[arr] (src) -- (parse);
        \draw[arr] (parse) -- (graph);
        \draw[arr] (graph) -- (seed);
        \draw[arr] (seed) -- (ppr);
        \draw[arr] (ppr) -- (deliver);
        \draw[arr] (deliver) -- (agent);

        % === Data flow arrows ===
        \draw[dataarr] (parse.east) -- (d1.west);
        \draw[dataarr] (graph.east) -- (d2.west);
        \draw[dataarr] (seed.east) -- (d3.west);
        \draw[dataarr] (ppr.east) -- (d4.west);

        % === Task description arrow ===
        \draw[arr, color=orange!60] (task) -- (seed);

        % === Phase grouping ===
        \begin{scope}[on background layer]
            \node[draw=blue!35, dashed, rounded corners=10pt, fill=blue!3,
                fit=(src)(parse)(graph)(d1)(d2),
                inner xsep=10pt, inner ysep=10pt,
                label={[font=\footnotesize\bfseries, color=blue!50!black]above left:
                        \raisebox{1pt}{\tikz\draw[blue!50, fill=blue!15, rounded corners=2pt] (0,0) rectangle (0.35,0.35);}\;%
                        Offline Indexing (5--30\,s)}] {};
            \node[draw=orange!40, dashed, rounded corners=10pt, fill=orange!2,
                fit=(task)(seed)(ppr)(deliver)(agent)(d3)(d4),
                inner xsep=10pt, inner ysep=10pt,
                label={[font=\footnotesize\bfseries, color=orange!60!black]above left:
                        \raisebox{1pt}{\tikz\draw[orange!50, fill=orange!15, rounded corners=2pt] (0,0) rectangle (0.35,0.35);}\;%
                        Online Retrieval ($<$\,250\,ms)}] {};
        \end{scope}
    \end{tikzpicture}
    \caption[CodeGraph pipeline architecture]{The five-stage CodeGraph pipeline. Stages~1--2 constitute \emph{offline indexing}: they parse the repository and construct the graph once per repository state (5--30\,s depending on repository size). Stages~3--5 constitute \emph{online retrieval}: they translate a task description into seeds, propagate relevance through the graph via PPR, and deliver ranked results to the agent (under 250\,ms). The right column shows the intermediate representation produced by each stage. The task description enters the pipeline at Stage~3 as the query input.}
    \label{fig:pipeline-architecture}
\end{figure}

The pipeline divides into two phases. The \emph{offline indexing} phase (Stages~1--2) runs once per repository state: it parses source code and constructs the dependency graph. The \emph{online retrieval} phase (Stages~3--5) runs per query: it translates a task description into seeds, propagates relevance through the graph, and delivers ranked results. Offline indexing requires 5--30\,seconds depending on repository size, while online retrieval completes in under 250\,ms.

The central design principle is \emph{signal amplification}. BM25 seed selection provides a noisy initial signal---an imperfect guess about which graph nodes are relevant to the task. Personalized PageRank then amplifies consistent structural patterns within that signal while attenuating noise. When multiple seeds independently point toward the same region of the graph, PPR concentrates probability there. When a seed is an outlier with no structural support, PPR dilutes its influence. This division of labor means that seed quality directly determines the retrieval ceiling---a finding confirmed empirically in \cref{ch:evaluation}.

To make this concrete, consider the task description ``Fix the SQL compilation error when using \texttt{QuerySet.annotate()}'' submitted against a Django repository with several hundred files. Entity extraction identifies \texttt{QuerySet}, \texttt{annotate}, and \texttt{SQLCompiler} as code identifiers and maps them to a handful of seed nodes. BM25 keyword search adds a few more nodes from files that mention SQL compilation. Some seeds are correct (within or adjacent to the relevant files); others are not. PPR then runs on the full graph: the correct seeds, which happen to reside in the same structural neighborhood---connected by \texttt{CONTAINS} and \texttt{CALLS} edges---reinforce each other, accumulating probability at their intersection. The incorrect seeds, which share no structural paths to the relevant region, dissipate their probability locally and are overwhelmed. The final ranking surfaces the structurally coherent neighborhood.

\Cref{tab:stage-complexity} summarizes the computational complexity of each stage. Offline indexing is dominated by Neo4j write operations; online retrieval is dominated by the PPR matrix-vector computation. In practice, offline indexing completes in 5--30\,s and online retrieval in under 250\,ms.

\begin{table}[H]
    \centering
    \caption[Per-stage computational complexity]{Computational complexity for each pipeline stage. $n$ = number of files; $m$ = average file size in tokens; $|V|$, $|E|$ = graph nodes and edges; $q$ = query length in tokens; $t$ = PPR convergence iterations (empirically 6--10); $k$ = top-$k$ result count. Offline stages run once per repository state; online stages run per query.}
    \label{tab:stage-complexity}
    \begin{tabularx}{\textwidth}{l l X l}
        \toprule
        \textbf{Stage} & \textbf{Phase} & \textbf{Dominant Operation} & \textbf{Complexity} \\
        \midrule
        Stage 1: Parsing          & Offline & tree-sitter CST traversal + entity extraction & $O(n \cdot m)$ \\
        \addlinespace[3pt]
        Stage 2: Graph Construction & Offline & Batched \texttt{UNWIND+MERGE} node/edge writes & $O(|V| + |E|)$ \\
        \addlinespace[3pt]
        IDF Reweighting           & Online  & Edge weight update across all edges & $O(|E|)$ \\
        \addlinespace[3pt]
        Stage 3: Seed Selection   & Online  & BM25 index query + entity name lookup & $O(q \cdot |V|)$ \\
        \addlinespace[3pt]
        Stage 4: PPR Ranking      & Online  & Sparse matrix-vector multiply (GDS) & $O(t \cdot |E|)$ \\
        \addlinespace[3pt]
        Stage 5: Context Delivery & Online  & Sort + token budget scan & $O(k \log k)$ \\
        \bottomrule
    \end{tabularx}
\end{table}

When errors occur across stages, the system degrades gracefully rather than failing hard. A file with a syntax error produces no entities and is absent from the graph---downstream stages simply never encounter it. Unresolved imports and ambiguous calls are dropped rather than guessed, preserving graph precision at the cost of recall. When seed selection produces an empty personalization vector---for instance, when a task description contains no recognizable identifiers and BM25 returns no results---the pipeline returns an empty result set and reports the failure explicitly. The guiding policy throughout is to prefer returning nothing over returning a misleading result.

Having established the pipeline's five-stage structure, the following sections examine each stage in detail, beginning with how raw source code is transformed into structured entities.


% =============================================================================
% 3.2 Source Code Parsing
% =============================================================================
\section{Source Code Parsing}
\label{sec:parsing}

Before CodeGraph can analyze repository structure, it must first understand what the code contains. A Python file is initially just text: the parser's job is to identify the meaningful units---classes, functions, methods---and the relationships between them, then hand that structure to the graph builder. This section describes how tree-sitter transforms raw source files into a typed entity model that serves as the direct input to graph construction.

The choice of entity granularity is consequential. A coarser model (files only) would lose function-level specificity; a finer model (individual statements) would produce an unmanageably dense graph. The four-entity model---files, classes, functions, methods---represents Python code at the natural level of importable and callable units, capturing the relationships that matter for dependency-driven retrieval without over-decomposing the structure.

\subsection{Parsing with tree-sitter}
\label{sec:parsing-treesitter}

CodeGraph relies on tree-sitter for all source code parsing. As introduced in \cref{ch:background}, tree-sitter is an incremental, error-tolerant parser that produces concrete syntax trees (CSTs) with millisecond latency per file. Its ability to handle partially-written or syntactically incorrect code is an important property when parsing real-world repositories that may contain syntax errors or experimental branches.

From each Python file, the parser extracts four types of code entities: files, classes, functions, and methods. Each entity receives a \emph{qualified name}---a string that uniquely identifies it within the repository. \Cref{tab:entity-types} summarizes the four entity types, their key properties, and the qualified name format for each.

\begin{table}[H]
    \centering
    \caption[Entity types extracted by the parser]{The four entity types extracted from Python source files. Each entity carries a qualified name that serves as a globally unique identifier within the repository graph. The double-colon separator (\texttt{::}) distinguishes the file path from the entity name, while the dot separator distinguishes class membership for methods.}
    \label{tab:entity-types}
    \begin{tabularx}{\textwidth}{l l X l}
        \toprule
        \textbf{Entity Type} & \textbf{Key Properties}                                   & \textbf{Description}                         & \textbf{Qualified Name}     \\
        \midrule
        File                 & \texttt{file\_path}, \texttt{line\_count}                 & A \texttt{.py} source file in the repository & \texttt{path/to/file.py}    \\
        \addlinespace[3pt]
        Class                & \texttt{name}, \texttt{file\_path}, \texttt{line\_number} & A class definition                           & \texttt{file::ClassName}    \\
        \addlinespace[3pt]
        Function             & \texttt{name}, \texttt{file\_path}, \texttt{signature}    & A module-level function definition           & \texttt{file::func\_name}   \\
        \addlinespace[3pt]
        Method               & \texttt{name}, \texttt{class\_name}, \texttt{signature}   & A function defined inside a class            & \texttt{file::Class.method} \\
        \bottomrule
    \end{tabularx}
\end{table}

The qualified name scheme ensures global uniqueness across the repository. Two files may each define a function named \texttt{process}, but \texttt{models/user.py::process} and \texttt{utils/http.py::process} are distinct nodes in the graph. This identity scheme serves as the primary node key throughout the system for deduplication, lookup, and cross-referencing.

\subsection{Import Resolution}
\label{sec:import-resolution}

Import statements declare cross-file dependencies, but tree-sitter provides only the raw import path as a string---it does not resolve which file on disk the path refers to. Import resolution converts these strings into concrete file references, which become \texttt{IMPORTS} edges in the graph.

The resolver handles three categories of imports. \emph{Absolute imports} (e.g., \texttt{from src.utils.helpers import validate\_email}) are resolved by converting the dotted path to a file system path relative to the project root and checking for a matching \texttt{.py} file or \texttt{\_\_init\_\_.py} package. \emph{Relative imports} (e.g., \texttt{from ..models import User}) are resolved by counting the leading dots to determine how many parent directories to traverse. \emph{Standard library and third-party imports} are identified using \texttt{sys.stdlib\_module\_names} (Python~3.10+) and a built-in names filter (\texttt{frozenset(dir(builtins))}), and are silently dropped rather than creating dangling nodes.

Unresolved imports---those referencing dynamic paths (\texttt{importlib.import\_module}), generated module names, or packages absent from the repository---are logged and skipped. This conservative strategy avoids introducing false edges into the graph. For well-structured Python repositories such as those in SWE-bench, the resolver successfully handles the large majority of import statements. The principal failure modes are third-party packages absent from the local repository tree (e.g., \texttt{requests}, \texttt{numpy}) and dynamic import patterns, which together account for the unresolved fraction. Standard library and built-in modules are correctly identified and filtered before resolution is attempted, so they do not inflate the failure count.

\subsection{Call Resolution}
\label{sec:call-resolution}

Call resolution converts function call expressions into \texttt{CALLS} edges. The goal is to determine, for each call site, which function or method definition the call refers to. This is more challenging than import resolution because Python is dynamically typed: the expression \texttt{user.save()} cannot be statically resolved without knowing the type of \texttt{user}.

The resolver employs a three-level priority scheme to match callee names against known definitions:

\begin{enumerate}[leftmargin=*, itemsep=4pt]
    \item \textbf{Imported names.} If the callee was imported via a successfully resolved import statement, the call is linked to the imported definition. This is the highest-confidence match.
    \item \textbf{Same-file definitions.} If the callee name matches a function or class defined in the same file, the call is linked to that local definition.
    \item \textbf{Unique global match.} If the callee name matches exactly one definition across the entire repository, the call is linked to that definition. When multiple definitions share the same name, the call is considered ambiguous.
\end{enumerate}

Ambiguous calls---where the callee name matches multiple definitions and cannot be disambiguated---are dropped. This is a deliberate design choice: a false \texttt{CALLS} edge connecting unrelated functions is more harmful to PPR accuracy than a missing edge, because PPR propagates probability along all edges indiscriminately. A false edge diverts probability mass into an irrelevant region of the graph, whereas a missing edge merely reduces the probability reaching a relevant node.

\subsection{Parser Limitations and Edge Cases}
\label{sec:parser-limitations}

Three edge cases merit explicit discussion because they affect a meaningful fraction of real-world repositories.

\textbf{Circular imports.} Python permits circular imports at the module level provided that symbols are not accessed at import time. The parser handles circular imports correctly by treating each file independently: import statements are recorded as dependency declarations regardless of whether the imported module has already been processed. The two-pass construction strategy (\cref{sec:two-pass}) then creates \texttt{IMPORTS} edges in the second pass when all files have been indexed, capturing both directions of a circular dependency without infinite recursion.

\textbf{Dynamic imports.} Calls to \texttt{importlib.import\_module()}, \texttt{\_\_import\_\_()}, or plugin-loading frameworks pass module names as runtime strings that are unresolvable by static analysis. These patterns are silently dropped. In heavily plugin-based codebases, this limitation leaves a substantial fraction of inter-module dependency edges absent from the graph, reducing the reachable set for queries that traverse those edges.

\textbf{Large files.} tree-sitter's incremental parser handles large files efficiently; parsing a 5,000-line file adds only marginally more latency than parsing five 1,000-line files, because CST construction is linear in file size. The practical bottleneck for large repositories is graph construction (Neo4j write batching), not parsing. Files larger than approximately 10,000 lines are unusual in well-maintained Python projects and were not encountered in the SWE-bench evaluation set.

\Cref{fig:parsing-extraction} illustrates the complete parsing pipeline from source code through entity extraction to resolved edges. With entities extracted and edges resolved, the next stage assembles them into a weighted property graph.

\begin{figure}[H]
    \centering
    \begin{tikzpicture}[
            % Code block
            codeblock/.style={draw, fill=gray!4, rounded corners=3pt, font=\ttfamily\scriptsize,
                    text width=5.8cm, align=left, inner sep=8pt,
                    line width=0.5pt, draw=gray!40},
            % Entity cards
            entity/.style={draw, rounded corners=4pt, minimum width=3.0cm,
                    minimum height=0.65cm, font=\small, align=center,
                    line width=0.5pt},
            % Arrows
            arr/.style={-{Stealth[length=5pt, width=4pt]}, thick, color=gray!60},
            % Section headers
            secheader/.style={font=\small\bfseries},
            node distance=0.35cm
        ]
        % === SOURCE CODE ===
        \node[secheader, color=gray!70] (clabel) {Source Code};
        \node[codeblock, below=0.15cm of clabel] (code) {%
            \textcolor{blue!70!black}{\textbf{from}} utils.crypto \textcolor{blue!70!black}{\textbf{import}} hash\_pw\\[3pt]
            \textcolor{blue!70!black}{\textbf{class}} \textcolor{green!50!black}{AuthService}\textcolor{gray}{:}\\[1pt]
            \quad\textcolor{blue!70!black}{\textbf{def}} \textcolor{orange!70!black}{authenticate}\textcolor{gray}{(self, pw):}\\[1pt]
            \quad\quad h = \textcolor{red!60!black}{hash\_pw}\textcolor{gray}{(pw)}\\[1pt]
            \quad\quad \textcolor{blue!70!black}{\textbf{return}} self.\textcolor{red!60!black}{verify}\textcolor{gray}{(h)}%
        };

        % === EXTRACTED ENTITIES ===
        \node[secheader=gray!70, right=3.2cm of clabel.east, anchor=west] (elabel) {Extracted Entities};

        \node[entity, fill=blue!8, below=0.15cm of elabel.south west, anchor=north west] (e1)
        {\footnotesize\texttt{File}\;\; \textcolor{gray!60}{|}\;\; \textcolor{blue!60}{\texttt{auth.py}}};
        \node[entity, fill=green!8, below=0.25cm of e1] (e2)
        {\footnotesize\texttt{Class}\;\; \textcolor{gray!60}{|}\;\; \textcolor{green!50!black}{\texttt{AuthService}}};
        \node[entity, fill=orange!10, below=0.25cm of e2] (e3)
        {\footnotesize\texttt{Method}\;\; \textcolor{gray!60}{|}\;\; \textcolor{orange!70!black}{\texttt{authenticate}}};
        \node[entity, fill=violet!6, below=0.25cm of e3] (e4)
        {\footnotesize\texttt{Import}\;\; \textcolor{gray!60}{|}\;\; \textcolor{violet!70}{\texttt{hash\_pw}}};
        \node[entity, fill=red!6, below=0.25cm of e4] (e5)
        {\footnotesize\texttt{Call}\;\; \textcolor{gray!60}{|}\;\; \textcolor{red!60!black}{\texttt{hash\_pw()}}};

        % === Arrows from code to entities ===
        \draw[arr, rounded corners=4pt] (code.east) -- ++(0.6,0) |- (e1.west);
        \draw[arr, rounded corners=4pt, color=green!40] ([yshift=-0.3cm]code.east) -- ++(0.35,0) |- (e2.west);
        \draw[arr, rounded corners=4pt, color=orange!40] ([yshift=-0.6cm]code.east) -- ++(0.2,0) |- (e3.west);

        % === RESOLVED EDGES ===
        \node[secheader=gray!70, below=0.8cm of e5.south west, anchor=north west] (rlabel) {Resolved Graph Edges};

        \node[draw, rounded corners=4pt, fill=gray!3, line width=0.4pt, draw=gray!30,
            inner sep=6pt, text width=6.5cm, font=\scriptsize,
            below=0.2cm of rlabel.south west, anchor=north west] (edges) {%
            \textcolor{blue!60}{\texttt{auth.py}} $\xrightarrow{\;\text{\scriptsize\sffamily CONTAINS}\;}$ \textcolor{green!50!black}{\texttt{AuthService}}\\[3pt]
            \textcolor{green!50!black}{\texttt{AuthService}} $\xrightarrow{\;\text{\scriptsize\sffamily CONTAINS}\;}$ \textcolor{orange!70!black}{\texttt{authenticate}}\\[3pt]
            \textcolor{blue!60}{\texttt{auth.py}} $\xrightarrow{\;\text{\scriptsize\sffamily IMPORTS}\;}$ \textcolor{blue!60}{\texttt{utils/crypto.py}}\\[3pt]
            \textcolor{orange!70!black}{\texttt{authenticate}} $\xrightarrow{\;\text{\scriptsize\sffamily CALLS}\;}$ \textcolor{red!60!black}{\texttt{hash\_pw}}%
        };

    \end{tikzpicture}
    \caption[Entity extraction and edge resolution from source code]{The parsing pipeline applied to a representative Python source file. \textbf{Left:} the source code, with color-coded identifiers corresponding to entity types. \textbf{Right, top:} the five entities extracted by tree-sitter---one File, one Class, one Method, one Import, and one Call. \textbf{Right, bottom:} the four graph edges resolved from these entities: two \texttt{CONTAINS} edges for structural containment, one \texttt{IMPORTS} edge from the import statement, and one \texttt{CALLS} edge from the function invocation.}
    \label{fig:parsing-extraction}
\end{figure}


% =============================================================================
% 3.3 Graph Construction
% =============================================================================
\section{Graph Construction}
\label{sec:graph-construction}

With entities extracted and references resolved, the second stage assembles them into a weighted property graph stored in Neo4j. This graph is the backbone of the entire system: it encodes every structural relationship in the repository---which file defines which class, which function calls which, which module imports which---and serves as the substrate on which PPR propagates relevance.

Two design choices at this stage have direct consequences for retrieval quality: the schema (which node and edge types to include) and the edge weight function (how to quantify the informativeness of each relationship). Both are addressed below.

\subsection{Node and Edge Schema}
\label{sec:graph-schema}

The graph uses four node labels (\texttt{File}, \texttt{Class}, \texttt{Function}, \texttt{Method}) corresponding directly to the entity types from \cref{sec:parsing}. Each node carries properties including \texttt{qualified\_name} (primary key), \texttt{name}, \texttt{file\_path}, and optional metadata such as \texttt{docstring}, \texttt{line\_number}, and \texttt{signature}.

Four edge types capture the structural relationships between code entities. \Cref{tab:edge-types} lists each edge type with its semantic meaning and endpoint types.

\begin{table}[H]
    \centering
    \caption[Edge types in the CodeGraph schema]{The four edge types in the code graph. All edges are created with a uniform base weight of~1.0, which is subsequently overwritten by IDF-based reweighting (\cref{sec:idf-reweighting}) before each PPR computation.}
    \label{tab:edge-types}
    \begin{tabularx}{\textwidth}{l l X c}
        \toprule
        \textbf{Edge Type}      & \textbf{Source $\to$ Target}              & \textbf{Semantic Meaning}                                                     & \textbf{Base Wt.} \\
        \midrule
        \texttt{CONTAINS}       & File $\to$ Class/Func; Class $\to$ Method & Structural containment: a file defines an entity, or a class defines a method & 1.0               \\
        \addlinespace[3pt]
        \texttt{CALLS}          & Func/Method $\to$ Func/Method             & A function or method invokes another                                          & 1.0               \\
        \addlinespace[3pt]
        \texttt{IMPORTS}        & File $\to$ File                           & One module imports from another                                               & 1.0               \\
        \addlinespace[3pt]
        \texttt{INHERITS\_FROM} & Class $\to$ Class                         & A class extends a parent class                                                & 1.0               \\
        \bottomrule
    \end{tabularx}
\end{table}

\Cref{fig:graph-schema} illustrates the graph schema as an entity-relationship diagram. Together, the four node types and four edge types capture the primary structural relationships in a Python codebase: containment (which file defines which class), invocation (which function calls which), dependency (which module imports which), and inheritance (which class extends which).

\begin{figure}[H]
    \centering
    \begin{tikzpicture}[
            % Node type boxes with properties
            nodetype/.style={draw, thick, rounded corners=5pt,
                    minimum width=2.6cm, minimum height=1.2cm,
                    font=\small\bfseries, align=center,
                    line width=0.7pt},
            % Edge styles with distinct patterns
            containsedge/.style={-{Stealth[length=5pt]}, thick, blue!60,
                    decoration={markings, mark=at position 0.5 with {\node[font=\sffamily\scriptsize\bfseries, fill=white, inner sep=1.5pt, text=blue!60] {CONTAINS};}},
                    postaction={decorate}},
            callsedge/.style={-{Stealth[length=5pt]}, thick, red!55,
                    decoration={markings, mark=at position 0.5 with {\node[font=\sffamily\scriptsize\bfseries, fill=white, inner sep=1.5pt, text=red!55] {CALLS};}},
                    postaction={decorate}},
            importsedge/.style={-{Stealth[length=5pt]}, thick, teal!70, dashed,
                    decoration={markings, mark=at position 0.5 with {\node[font=\sffamily\scriptsize\bfseries, fill=white, inner sep=1.5pt, text=teal!70] {IMPORTS};}},
                    postaction={decorate}},
            inheritsedge/.style={-{Stealth[length=5pt]}, thick, violet!65, densely dotted,
                    decoration={markings, mark=at position 0.5 with {\node[font=\sffamily\scriptsize\bfseries, fill=white, inner sep=1.5pt, text=violet!65] {INHERITS};}},
                    postaction={decorate}},
            node distance=3.0cm and 4.0cm
        ]
        % Nodes with properties shown
        \node[nodetype, fill=blue!12, draw=blue!40] (file) {\textcolor{blue!60}{File}\\[-1pt]\scriptsize\textcolor{gray}{file\_path, line\_count}};
        \node[nodetype, fill=green!10, draw=green!40, right=4.0cm of file] (class) {\textcolor{green!50!black}{Class}\\[-1pt]\scriptsize\textcolor{gray}{name, bases, docstring}};
        \node[nodetype, fill=orange!12, draw=orange!50, below=3.0cm of file] (func) {\textcolor{orange!70!black}{Function}\\[-1pt]\scriptsize\textcolor{gray}{name, signature}};
        \node[nodetype, fill=purple!10, draw=purple!40, below=3.0cm of class] (method) {\textcolor{purple!60!black}{Method}\\[-1pt]\scriptsize\textcolor{gray}{name, class\_name, sig}};

        % CONTAINS edges (solid blue)
        \draw[containsedge] (file) -- (class);
        \draw[containsedge] (file) -- (func);
        \draw[containsedge] (class) -- (method);

        % CALLS edges (solid red, curved)
        \draw[callsedge] (func) -- (method);
        \draw[callsedge, bend left=18] (method.west) to (func.east);

        % IMPORTS edge (dashed teal, self-loop)
        \draw[importsedge, loop above, looseness=5, out=115, in=65, distance=1.4cm]
        (file) to (file);

        % INHERITS_FROM edge (dotted violet, self-loop)
        \draw[inheritsedge, loop above, looseness=5, out=115, in=65, distance=1.4cm]
        (class) to (class);

        % Legend box with visual line samples
        \node[draw=gray!30, rounded corners=4pt, fill=gray!3, font=\scriptsize,
            text width=6.5cm, align=left, inner sep=7pt,
            below=0.3cm of $(func.south)!0.5!(method.south)$, anchor=north] (legend) {%
            \textbf{Edge Legend}\\[3pt]
            \textcolor{blue!60}{\rule[2pt]{0.8cm}{1.5pt}} \textsf{CONTAINS} --- structural containment (file$\to$entity, class$\to$method)\\[2pt]
            \textcolor{red!55}{\rule[2pt]{0.8cm}{1.5pt}} \textsf{CALLS} --- function/method invocation\\[2pt]
            {\color{teal!70}\rule[2pt]{0.15cm}{1.5pt}\hspace{2pt}\rule[2pt]{0.15cm}{1.5pt}\hspace{2pt}\rule[2pt]{0.15cm}{1.5pt}\hspace{2pt}\rule[2pt]{0.15cm}{1.5pt}} \textsf{IMPORTS} --- module-level dependency\\[2pt]
            {\color{violet!65}\rule[2pt]{0.08cm}{1.5pt}\hspace{1.5pt}\rule[2pt]{0.08cm}{1.5pt}\hspace{1.5pt}\rule[2pt]{0.08cm}{1.5pt}\hspace{1.5pt}\rule[2pt]{0.08cm}{1.5pt}\hspace{1.5pt}\rule[2pt]{0.08cm}{1.5pt}\hspace{1.5pt}\rule[2pt]{0.08cm}{1.5pt}\hspace{1.5pt}\rule[2pt]{0.08cm}{1.5pt}} \textsf{INHERITS\_FROM} --- class inheritance
        };
    \end{tikzpicture}
    \caption[Graph schema diagram]{The CodeGraph property graph schema. Four node types (colored boxes) represent code entities; four edge types (distinguished by color and line style) capture structural relationships. Each node type displays its key properties in gray. Self-referencing arrows on \texttt{File} and \texttt{Class} indicate that \texttt{IMPORTS} and \texttt{INHERITS\_FROM} edges connect nodes of the same type (e.g., one file imports another; one class inherits from another).}
    \label{fig:graph-schema}
\end{figure}


\subsection{Two-Pass Construction}
\label{sec:two-pass}

Graph construction proceeds in two passes. The first pass creates all nodes and intra-file \texttt{CONTAINS} edges. The second pass creates all cross-file edges (\texttt{CALLS}, \texttt{IMPORTS}, \texttt{INHERITS\_FROM}) after every definition in the repository has been indexed.

The two-pass strategy is necessary because of forward references. A call in file~A may target a function defined in file~B, but if file~B has not yet been parsed, the callee definition does not exist in the graph. By deferring cross-file edge creation to the second pass, all definitions are available for resolution regardless of file processing order. A single-pass strategy using topological file ordering was considered but rejected: import cycles---which make topological ordering undefined---are present in a substantial fraction of real-world Python repositories. The two-pass approach handles cycles correctly without additional complexity.

The procedure is as follows.

\noindent\textbf{Pass 1 --- intra-file construction:}
\begin{enumerate}[leftmargin=*, itemsep=3pt]
    \item Parse each Python file with tree-sitter to produce a concrete syntax tree.
    \item Extract all entities: \texttt{File}, \texttt{Class}, \texttt{Function}, and \texttt{Method} nodes.
    \item Write entity nodes to Neo4j using batched \texttt{UNWIND+MERGE} operations (idempotent).
    \item Create \texttt{CONTAINS} edges: File $\to$ Class/Function, and Class $\to$ Method.
    \item Record all import statements and call expressions for resolution in Pass~2.
\end{enumerate}

\noindent\textbf{Pass 2 --- cross-file linking:}
\begin{enumerate}[leftmargin=*, itemsep=3pt]
    \item For each recorded import, resolve the path to a target file and create an \texttt{IMPORTS} edge on success; log and skip failures.
    \item For each recorded function call, match the callee using the three-level scheme from \cref{sec:call-resolution} and create a \texttt{CALLS} edge for unambiguous matches only.
    \item For each class with declared base classes, resolve the base class names to \texttt{Class} nodes and create \texttt{INHERITS\_FROM} edges.
\end{enumerate}

All writes use \texttt{UNWIND+MERGE}, guaranteeing idempotency: re-parsing a file after an edit produces no duplicate nodes or edges. Uniqueness constraints on the \texttt{qualified\_name} property at the database level provide an additional safety guarantee.


\subsection{IDF-Based Edge Reweighting}
\label{sec:idf-reweighting}

The uniform base weights assigned during construction do not reflect the varying informativeness of different edges. A \texttt{CALLS} edge to a ubiquitous utility function invoked by many callers (e.g., \texttt{validate\_input}, \texttt{get\_config}) conveys less structural signal about relevance than an edge to a specialized function invoked by only a few callers. To capture this asymmetry, CodeGraph applies inverse document frequency (IDF) reweighting to edge weights before each PPR computation.

The IDF weight for an edge targeting node~$v$ is:

\begin{equation}
    w(e) = \frac{1}{\log_2\!\bigl(\operatorname{in\_degree}(v) + 2\bigr)}
    \label{eq:idf-weight}
\end{equation}

\noindent where $\operatorname{in\_degree}(v)$ is the number of incoming edges to the target node. The base-2 logarithm provides an interpretable doubling semantics: each doubling of a node's in-degree reduces the edge weight by exactly one $\log_2$ unit. A node with in-degree~2 receives weight $1/\log_2(4) = 0.50$; a node with in-degree~14 receives weight $1/\log_2(16) = 0.25$---half the weight for eight times the connectivity. The $+2$ offset serves two purposes: it ensures that zero-in-degree nodes receive a finite, positive weight ($1/\log_2(2) = 1.0$), and it prevents division by zero. Alternative formulations were evaluated on the benchmark: replacing the denominator with $\log_{10}(\text{in\_deg} + 1)$ produced a 2.3\,pp drop in Recall@10, and the simpler inverse-linear formula $1/(\text{in\_deg} + 1)$ produced an 8.1\,pp drop, suggesting that logarithmic compression is essential for preserving useful signal from moderately-connected nodes. \Cref{tab:idf-examples} illustrates how IDF weights vary with in-degree.

\begin{table}[H]
    \centering
    \caption[IDF weight examples]{IDF edge weight as a function of the target node's in-degree. Highly connected nodes (utility functions, common initializers) receive lower weights, reflecting their lower structural informativeness for retrieval purposes.}
    \label{tab:idf-examples}
    \begin{tabular}{l r r l}
        \toprule
        \textbf{Target Entity Example}             & \textbf{In-degree} & \textbf{IDF Weight} & \textbf{Interpretation}          \\
        \midrule
        Specialized helper function                & 3                  & 0.42                & High: few callers, strong signal \\
        Module-level utility                       & 15                 & 0.24                & Moderate                         \\
        Common initializer (\texttt{\_\_init\_\_}) & 100                & 0.15                & Low: ubiquitous, weak signal     \\
        \bottomrule
    \end{tabular}
\end{table}

CodeGraph applies IDF reweighting at retrieval time, not at construction time. This ensures that edge weights always reflect the current state of the graph, which may shift when files are re-parsed and new edges are added. The reweighting runs on every retrieval call as part of the \texttt{ensure\_graph\_ready()} step, incurring approximately 50\,ms overhead---marginal relative to the 250\,ms total query latency.


\subsection{GDS Projection}
\label{sec:gds-projection}

The Neo4j Graph Data Science (GDS) library performs PPR computation on an in-memory graph projection rather than on the disk-backed storage graph. Before PPR can run, the relevant subgraph must be projected into GDS memory.

The projection includes all four node types and all four edge types. Edges are projected as \textbf{undirected}, allowing probability to flow in both directions during the random walk. This is a deliberate choice: an AI agent fixing a function needs to see both what the function depends on (downstream---following call edges forward) and what depends on the function (upstream---following call edges backward). The \texttt{weight} property set by IDF reweighting is included as the relationship weight in the projection.

The projection is dropped and recreated on every retrieval call to ensure consistency with the current IDF weights. For a graph with 50,000 nodes and 200,000 edges, this operation takes approximately 50\,ms and requires roughly 5\,MB of memory---both negligible at the scale of typical development repositories.


% =============================================================================
% 3.4 Seed Selection
% =============================================================================
\section{Seed Selection}
\label{sec:seed-selection}

The graph now exists in Neo4j, encoding every structural dependency in the repository. PPR can propagate relevance through this graph---but it needs a starting point. A pure graph algorithm has no way to distinguish a relevant file from an irrelevant one without some initial signal derived from the task description.

Seed selection answers the question: \emph{where in the graph should the random walk begin?} It translates a natural language task description into a \emph{personalization vector}---a probability distribution over graph nodes that specifies which nodes serve as PPR restart points. The seeds do not need to be perfect; PPR's structural propagation tolerates a fraction of incorrect seeds. But seed quality sets a hard ceiling: when no seed falls in the relevant region of the graph, PPR has no path to the answer regardless of how many iterations it runs. As shown in \cref{ch:evaluation}, seed quality accounts for the dominant share of retrieval improvement across the system's development history.

CodeGraph constructs the personalization vector from three independent signals, each capturing a different dimension of evidence. \Cref{fig:seed-selection} illustrates the signal flow. The default weights (entity match~0.6, BM25~0.3, current file~0.1) reflect the relative reliability of each source: explicit identifier mentions provide the strongest direct evidence, keyword search provides complementary coverage when identifiers are absent, and the current file provides only a weak spatial hint. The specific weights are empirically motivated; \cref{ch:evaluation} examines how seed quality as a whole affects retrieval performance.

\begin{figure}[H]
    \centering
    \begin{tikzpicture}[
            % Input boxes
            inputbox/.style={draw, rounded corners=4pt, fill=gray!6, minimum height=0.95cm,
                    text width=2.4cm, align=center, font=\small,
                    line width=0.5pt, draw=gray!40},
            % Signal processing blocks
            signal/.style={draw, rounded corners=4pt, minimum height=0.85cm,
                    text width=2.5cm, align=center, font=\scriptsize,
                    line width=0.5pt},
            % Merge node
            merge/.style={draw, rounded corners=6pt, fill=yellow!12, draw=yellow!50,
                    minimum width=3.2cm, minimum height=0.95cm,
                    font=\small\bfseries, align=center, line width=0.6pt},
            % Output
            output/.style={draw, rounded corners=5pt, fill=orange!12, draw=orange!50,
                    minimum width=3.0cm, minimum height=0.9cm,
                    font=\small\bfseries, align=center, line width=0.6pt},
            % Arrows
            arr/.style={-{Stealth[length=5pt, width=4pt]}, thick, color=gray!60},
            % Weight labels
            wt/.style={font=\scriptsize\bfseries, fill=white, inner sep=1.5pt, rounded corners=1pt,
                    draw=gray!20},
            node distance=0.55cm and 0.7cm
        ]
        % === INPUTS ===
        \node[inputbox] (task) {Task\\Description};
        \node[inputbox, below=2.8cm of task] (file) {Current File\\{\scriptsize(optional)}};

        % === THREE SIGNAL STREAMS ===
        % Stream 1: Entity extraction
        \node[signal, fill=blue!10, right=1.8cm of task] (regex) {Entity Name\\Extraction\\{\tiny(regex patterns)}};
        \node[signal, fill=blue!6, below=0.55cm of regex] (lookup) {Graph Lookup\\+ Inverse-freq\\weighting};

        % Stream 2: BM25
        \node[signal, fill=green!10, right=0.7cm of regex] (bm25tok) {Compound\\Tokenizer\\{\tiny(CamelCase split)}};
        \node[signal, fill=green!6, below=0.55cm of bm25tok] (bm25idx) {BM25 Index\\$\to$ Top-$N$\\scored results};

        % Stream 3: Current file
        \node[signal, fill=purple!8, right=0.7cm of bm25tok] (filsig) {File Entity\\Lookup\\{\tiny(all entities in file)}};
        \node[signal, fill=purple!5, below=0.55cm of filsig] (fileent) {Weak\\proximity\\hint};

        % === MERGE ===
        \node[merge, below=1.0cm of bm25idx] (mergenode) {Merge \& Normalize\\{\scriptsize$\sum_i s_i = 1.0$}};

        % === OUTPUT ===
        \node[output, below=0.6cm of mergenode] (pvec) {Personalization Vector\;\;$\vec{s}$};

        % === INPUT ARROWS ===
        \draw[arr] (task.east) -- (regex.west);
        \draw[arr] (task.east) -- ++(0.6,0) |- (bm25tok.west);
        \draw[arr] (file.east) -- ++(3.5,0) |- (filsig.south west);

        % === STREAM ARROWS ===
        \draw[arr, color=blue!50] (regex) -- (lookup);
        \draw[arr, color=green!50!black] (bm25tok) -- (bm25idx);
        \draw[arr, color=purple!40] (filsig) -- (fileent);

        % === MERGE ARROWS with weights ===
        \draw[arr, color=blue!50] (lookup.south) -- node[wt, left=1pt] {\textcolor{blue!60}{0.6}} (mergenode.north west);
        \draw[arr, color=green!50!black] (bm25idx.south) -- node[wt] {\textcolor{green!50!black}{0.3}} (mergenode.north);
        \draw[arr, color=purple!40] (fileent.south) -- node[wt, right=1pt] {\textcolor{purple!50}{0.1}} (mergenode.north east);

        % === OUTPUT ARROW ===
        \draw[arr, color=orange!60, line width=1.2pt] (mergenode) -- (pvec);

        % === STREAM LABELS ===
        \node[font=\tiny\bfseries, color=blue!60, above=0.05cm of regex] {HIGH CONFIDENCE};
        \node[font=\tiny\bfseries, color=green!50!black, above=0.05cm of bm25tok] {MEDIUM};
        \node[font=\tiny\bfseries, color=purple!50, above=0.05cm of filsig] {LOW};

    \end{tikzpicture}
    \caption[Seed selection signal flow]{Seed selection combines three signals with configurable weights. \textbf{Entity extraction} (weight~0.6, blue) uses regex patterns to identify code identifiers in the task text and matches them against graph nodes, with inverse-frequency weighting to penalize ambiguous names. \textbf{BM25 keyword search} (weight~0.3, green) applies compound tokenization to split identifiers at camelCase/underscore boundaries, then ranks entities by textual relevance. \textbf{Current file} (weight~0.1, purple) adds all entities from the agent's active file as a weak proximity hint. The three signals are merged by summing weights per node and normalizing to produce a valid probability distribution~$\vec{s}$.}
    \label{fig:seed-selection}
\end{figure}


\subsection{Entity Match Signal}
\label{sec:entity-match}

The entity match signal extracts likely code identifiers from the task description using a set of regex patterns, then matches them against nodes in the graph. The patterns target naming conventions common in Python: multi-word CamelCase identifiers (\texttt{AuthService}, \texttt{TokenExpiredError}), backtick-quoted identifiers from GitHub Markdown (\texttt{`models.QuerySet`}), dotted paths (\texttt{sql.compiler}), and snake\_case identifiers (\texttt{handle\_subquery}).

Each extracted name is looked up in the graph. When a name matches multiple nodes---for example, \texttt{update} might match dozens of methods across a Django codebase---an \emph{inverse-frequency weighting} correction is applied. The weight assigned to each matching node is scaled by $1 / \log_2(n + 1)$, where $n$ is the number of matches. This prevents ambiguous names from dominating the seed distribution: the name \texttt{update} matching 30~methods contributes far less total weight than the name \texttt{SQLUpdateCompiler} matching a single class.

The entity match signal receives the highest default weight (0.6) because explicit identifier mentions in a bug report---such as ``fix the bug in \texttt{SQLCompiler}''---serve as direct pointers to specific graph nodes, providing the strongest available evidence of relevance.


\subsection{BM25 Signal}
\label{sec:bm25-signal}

The BM25 signal performs keyword search over a corpus of text documents, one per graph entity. Each document concatenates several text sources: the entity's short name (e.g., \texttt{SQLCompiler}), its qualified name (e.g., \texttt{django/db/models/sql/compiler.py::SQLCompiler}), the file path tokenized as path components (e.g., \texttt{django db models sql compiler}), and the first 200 characters of the docstring.

A critical preprocessing step is \emph{compound tokenization}: splitting identifier names at camelCase boundaries and underscores before indexing. Without this splitting, \texttt{SQLCompiler} and the words ``SQL'' and ``compiler'' are three distinct, unrelated tokens in the BM25 vocabulary, and a query mentioning ``SQL compilation'' fails to match the class. Compound tokenization transforms \texttt{SQLCompiler} into the tokens \texttt{sql} and \texttt{compiler}, enabling partial matching. This preprocessing alone improves the BM25-only baseline by +8.3\,pp Recall@10 (from 49.7\% to 58.0\%), as quantified in \cref{ch:evaluation}.

The BM25 search returns the top-10 results, each scored proportionally by BM25 relevance. These are added to the personalization vector with a default signal weight of~0.3.


\subsection{Current File Signal}
\label{sec:current-file-signal}

When the calling agent provides the path to the file currently being edited, all entities defined in that file are added to the personalization vector with a low weight (default~0.1). This signal acts as a weak proximity hint: the currently active file is often in the neighborhood of the relevant code, even if it is not the relevant file itself.

The current file signal is intentionally weighted low because where the cursor is does not reliably correspond to what the developer cares about. The developer may have navigated to the current file during an exploratory search, and the file may bear no connection to the actual task.


\subsection{Personalization Vector Construction}
\label{sec:pvec-construction}

The three signals are merged into a single personalization vector $\vec{s} \in \mathbb{R}^{|V|}$ by summing the weights for each node across all signals. If a node appears in both the entity match results and the BM25 results, its weights are added. The merged vector is then normalized to satisfy $\|\vec{s}\|_1 = 1$, producing a valid probability distribution over graph nodes. This vector is passed directly to the GDS PPR implementation as the restart distribution.

A worked example illustrates the full signal flow. Given the task description ``Fix the authentication failure when \texttt{TokenExpiredError} is raised'', entity extraction identifies the identifier \texttt{TokenExpiredError} and matches it to a single node: \texttt{auth/exceptions.py::TokenExpiredError}. BM25 keyword search on the compound-tokenized index matches several additional entities: \texttt{auth/middleware.py::authenticate}, \texttt{auth/tokens.py::verify\_token}, and \texttt{auth/views.py::LoginView}. No current file is provided. The resulting personalization vector places 0.6 of its mass on the single entity match, distributes 0.3 proportionally across the four BM25 results, and normalizes---concentrating the seed distribution on a tight neighborhood in the authentication subsystem. PPR then propagates from this neighborhood through the graph structure.

When a task description contains no recognizable identifiers---for example, ``The application crashes on startup''---entity extraction yields nothing and BM25 may only weakly match common tokens like ``application''. In such cases, the BM25 signal dominates by default after normalization, and the personalization vector reflects only shallow keyword evidence. These description-impoverished instances contribute disproportionately to zero-recall cases in evaluation, as discussed in \cref{ch:evaluation}. This pattern identifies a fundamental limitation of static text matching: when the task description does not mention the code by name, no retrieval signal can overcome the absence of identifying information.

With the personalization vector constructed, the pipeline hands it to Personalized PageRank to propagate relevance through the graph structure.


% =============================================================================
% 3.5 Personalized PageRank Configuration
% =============================================================================
\section{Personalized PageRank Configuration}
\label{sec:ppr-config}

With seeds selected, the pipeline invokes Personalized PageRank to propagate relevance through the graph. As introduced in \cref{ch:background}, PPR simulates a random walk that follows graph edges with probability $d$ (the \emph{damping factor}) and teleports back to the seed nodes with probability $1-d$. The stationary distribution of this walk assigns each node a relevance score that reflects both its structural centrality and its proximity to the seeds.

PPR exposes three parameters with meaningful consequences for retrieval quality: the restart mode (how probability is distributed across seeds), the damping factor (how far the walk ventures from seeds), and the top-$k$ cutoff (how many results to return). Each parameter was set through controlled benchmark evaluation rather than theoretical defaults. This section presents the rationale behind each choice and the evidence supporting it.

\subsection{Uniform vs.\ Weighted Restart Mode}
\label{sec:uniform-vs-weighted}

When PPR restarts from multiple seed nodes, a natural question arises: should seeds with higher confidence---say, an exact entity name match---receive more restart probability than seeds derived from weaker evidence like a partial BM25 match? CodeGraph implements two strategies that answer this differently.

\paragraph{Weighted restart mode.} Each seed receives restart probability proportional to its weight in the personalization vector $\vec{s}$. Seeds with higher BM25 scores or stronger entity matches receive more restart probability. This mode computes:
\begin{equation}
    \text{PPR}_{\text{weighted}} = \sum_i w_i \cdot \text{PPR}(\text{seed}_i)
    \label{eq:weighted-ppr}
\end{equation}
where $w_i$ is seed $i$'s normalized weight. The combination is mathematically exact because PPR is linear in the restart distribution.

\paragraph{Uniform restart mode (default).} All seeds receive equal restart probability regardless of their weights. This mode executes a single GDS \texttt{pageRank.stream()} call with all seed node IDs passed as \texttt{sourceNodes}---one graph traversal, not $|\text{seeds}|$.

\Cref{tab:restart-modes} compares the two modes across two evaluation iterations with progressively improved seed quality.

\begin{table}[H]
    \centering
    \caption[Uniform vs.\ weighted restart mode comparison]{Recall@10 and zero-recall instance counts for uniform vs.\ weighted restart across two evaluation iterations. Uniform mode consistently outperforms weighted mode; the gap narrows as seed quality improves but does not disappear. Zero-recall counts reflect how often no relevant file appeared in the top-10 results ($n = 300$ benchmark instances per row).}
    \label{tab:restart-modes}
    \begin{tabularx}{\textwidth}{l X r r r r}
        \toprule
        \textbf{Iteration} & \textbf{Seed Quality} & \textbf{Weighted R@10} & \textbf{Uniform R@10} & \textbf{$\Delta$} & \textbf{Zero-recall (uniform)} \\
        \midrule
        1 & Basic entity match & 42.0\% & 60.7\% & +18.7\,pp & 117 \\
        \addlinespace[3pt]
        2 & +BM25 +compound tokens & 66.7\% & 72.3\% & +5.6\,pp & 83 \\
        \bottomrule
    \end{tabularx}
\end{table}

\noindent The direction of the advantage is consistent across all 300 benchmark instances in both iterations, suggesting a structural cause rather than an incidental one.

\textbf{Why does uniform mode win?} The explanation centers on overconfident seed weights. When entity extraction picks up an ambiguous name---for example, \texttt{fit} matching 30~scikit-learn estimators---the entity match signal (weight~0.6) distributes substantial restart probability across incorrect nodes. BM25 seeds (weight~0.3) cannot compensate because their contribution is diluted in the weighted sum. Uniform mode avoids this failure by giving every seed---including the BM25 seeds that often locate the correct file neighborhood---equal restart probability. From that point, the graph structure itself arbitrates: the coherent seeds, which share structural paths through shared dependencies, receive compounding PPR mass at their intersection. The isolated incorrect seeds, which lack structural support, dissipate probability locally and are overtaken.

This finding reflects a general principle in seed-based graph retrieval: \emph{when seed quality is imperfect, seed diversity matters more than seed confidence.} Trusting the graph to identify structural consistency is more robust than trusting the seed weighting to identify correctness.

\subsection{Damping Factor}
\label{sec:damping-factor}

The damping factor $d$ controls the trade-off between following graph edges and restarting from seeds. CodeGraph uses $d = 0.70$ (restart probability $1 - d = 0.30$), compared to the standard default of $d = 0.85$.

A lower damping factor keeps the random walk closer to the seed nodes on each step. With improved seed quality (after the BM25 and entity extraction improvements from \cref{sec:seed-selection}), staying closer to the seeds improves early-rank quality: Recall@5 increased by +3\,pp (from 60.7\% to 63.7\%), while Recall@10 remained unchanged at 72.3\%. The lower damping factor also has a beneficial side effect on convergence speed: with $d = 0.70$, the random walk restarts more frequently, reducing the expected walk length and therefore the number of matrix-vector multiply iterations required to reach the $10^{-6}$ tolerance threshold. Empirically, PPR with $d = 0.70$ converges in approximately 6--10 iterations on the benchmark graphs.

The Recall@10 ceiling is unaffected because it is determined by which files receive \emph{any} nonzero PPR probability---a property governed by graph connectivity, not by the damping factor. The damping factor affects only the \emph{ranking} within the reachable set. In graph-theoretic terms, the expected random walk length under damping factor $d$ is $1/(1-d)$ hops: at $d = 0.70$ this is $\approx 3.3$ hops; at $d = 0.85$ it is $\approx 6.7$ hops. The longer walk at $d = 0.85$ reaches marginally more nodes, but the difference in reachable set cardinality is small relative to repository size, which explains why Recall@10 (a reachability metric at this scale) is insensitive to the change while Recall@5 (a ranking-sensitive metric) is not.

\subsection{Top-\textit{k} Configuration}
\label{sec:topk}

The system returns the top-$k$ files by PPR score, with a default of $k = 30$. Increasing $k$ from 20 to 30 resolved 5~additional benchmark instances whose gold files were ranked at positions~21--30, yielding a direct +1.7\,pp improvement in Recall@10 at zero algorithmic cost. The marginal cost is small: returning 10~additional file paths adds negligible latency, and the token budget mechanism in post-processing (\cref{sec:post-processing}) independently limits how much source code is actually served to the agent.

With PPR configuration established, the ranked list of entities moves into the post-processing stage, where graph-level scores are translated into agent-ready source code.


% =============================================================================
% 3.6 Post-Processing and Context Delivery
% =============================================================================
\section{Post-Processing and Context Delivery}
\label{sec:post-processing}

PPR returns a ranked list of individual entities---functions, methods, classes---but an AI agent needs source code, not node identifiers. Post-processing bridges this gap through four steps: deduplication to avoid redundant file representations, source code extraction at entity-level line ranges, token budget enforcement to respect context window limits, and structured formatting for the MCP response. Each step addresses a distinct constraint between the graph-level result and the agent-level artifact.

\paragraph{File-level deduplication.} PPR scores individual entities (functions, methods, classes), but the agent ultimately needs file-level context. When two or more sub-entities from the same file appear in the top-$k$ results, the corresponding \texttt{File} node is dropped to avoid consuming token budget on a file already represented by its constituent entities. The threshold of two was chosen empirically: a threshold of one would collapse all results to file granularity, discarding the entity-level ranking signal; a threshold of three would allow redundant whole-file representations to inflate token usage when entity-level results already cover the relevant content.

\paragraph{Source code reading.} For each ranked entity, the system reads the relevant source code from disk using the entity's \texttt{file\_path}, \texttt{line\_number}, and \texttt{end\_line} properties. Entity-level line ranges ensure that only the relevant portion of a file is included rather than the entire file.

\paragraph{Token budget enforcement.} The total token count of the returned source code is bounded by a configurable budget (default: 15{,}000 tokens, measured with the \texttt{cl100k\_base} tokenizer via \texttt{tiktoken}). Results are added in rank order until the budget is exhausted. At least one result is always returned even if it exceeds the budget, ensuring that the agent receives some context for every query. The 15{,}000-token default was chosen to balance context richness with practical constraints: empirically, most SWE-bench tasks require between 5 and 15 files of source code for sufficient context, and this budget typically accommodates 8--12 files depending on file length. Larger budgets marginally improve recall by including additional ranked results, but at the cost of increasing the context window pressure on the downstream language model---a trade-off that can be tuned via configuration without changes to the pipeline.

\paragraph{Output format.} Each result comprises eight fields: entity name, entity type, qualified name, file path, line range, relevance score, source code text, and token count. This structured output allows the MCP server to format results into a tool response that the agent can parse and act on.


% =============================================================================
% 3.7 MCP Server Integration
% =============================================================================
\section{MCP Server Integration}
\label{sec:mcp-integration}

The preceding sections described the retrieval pipeline as a self-contained computation. This section explains how that pipeline is exposed to AI agents through the Model Context Protocol (MCP)---an open standard for agent tool use that defines a message-passing interface for invoking external capabilities via structured tool calls and receiving structured results.

CodeGraph implements an MCP server that wraps the retrieval pipeline as a set of callable tools. The server is stateless per request---no session state is maintained between tool calls---while the Neo4j graph persists across calls, eliminating the need to rebuild the graph on each invocation.

\subsection{Tool Interface}
\label{sec:tool-interface}

\Cref{tab:mcp-tools} lists the five core tools exposed by the MCP server. The primary entry point is \texttt{get\_relevant\_context}, which accepts a task description and runs the full retrieval pipeline (seed selection $\to$ PPR $\to$ post-processing). The remaining tools provide complementary capabilities for dependency exploration, graph inspection, and ad-hoc querying.

\begin{table}[H]
    \centering
    \caption[MCP tool reference]{The five core tools exposed by the CodeGraph MCP server. Each tool accepts structured parameters and returns a JSON result. The primary tool (\texttt{get\_relevant\_context}) orchestrates the full retrieval pipeline; the remaining tools support targeted exploration and debugging during agent-driven development sessions.}
    \label{tab:mcp-tools}
    \begin{tabularx}{\textwidth}{l X X}
        \toprule
        \textbf{Tool}                   & \textbf{Parameters}                                            & \textbf{Use Case}                                                                     \\
        \midrule
        \texttt{get\_relevant\_context} & Task description, optional entity hints, top-$k$, token budget & Primary entry point: runs the full pipeline and returns ranked files with source code \\
        \addlinespace[4pt]
        \texttt{query\_dependencies}    & Qualified entity name, direction (up/downstream), depth        & Explores the dependency neighborhood of a specific entity to assess blast radius      \\
        \addlinespace[4pt]
        \texttt{get\_graph\_stats}      & (none)                                                         & Returns graph summary: node/edge counts by type, connected file count                 \\
        \addlinespace[4pt]
        \texttt{find\_dead\_code}       & (none)                                                         & Identifies entities with no incoming \texttt{CALLS} or \texttt{IMPORTS} edges         \\
        \addlinespace[4pt]
        \texttt{execute\_cypher\_query} & Cypher query string                                            & Raw graph query access for agents requiring custom traversal logic                    \\
        \bottomrule
    \end{tabularx}
\end{table}

In addition to the five core tools, the server exposes five supplementary tools (\texttt{find\_callers}, \texttt{find\_callees}, \texttt{visualize\_query}, \texttt{analyze\_complexity}, \texttt{watch\_directory}) for more targeted analysis during interactive development sessions.


\subsection{Agent Workflow}
\label{sec:agent-workflow}

\Cref{fig:agent-workflow} shows the typical interaction between an AI agent and the CodeGraph MCP server. The agent receives a task description (e.g., a bug report), calls \texttt{get\_relevant\_context} with the task text, and receives ranked source code in return. It may then optionally call \texttt{query\_dependencies} to explore the neighborhood of suspicious entities before generating a solution.

\begin{figure}[H]
    \centering
    \begin{tikzpicture}[
            % Actor headers
            actor/.style={draw, rounded corners=5pt, minimum width=2.8cm,
                    minimum height=1.0cm, font=\small\bfseries, align=center,
                    line width=0.6pt},
            % Message arrows
            msgarr/.style={-{Stealth[length=5pt, width=4pt]}, thick},
            msgarrdash/.style={-{Stealth[length=5pt, width=4pt]}, thick, dashed},
            % Message labels
            msg/.style={font=\scriptsize, fill=white, inner sep=2pt, rounded corners=1pt},
            % Step numbers
            stepnum/.style={font=\scriptsize\bfseries, circle,
                    minimum size=14pt, inner sep=0pt},
            % Timeline
            tl/.style={thick, gray!40},
            node distance=0.5cm
        ]
        % === ACTORS ===
        \node[actor, fill=gray!8, draw=gray!40] (agent) {AI Agent};
        \node[actor, fill=purple!8, draw=purple!35, right=3.8cm of agent] (mcp) {MCP Server};
        \node[actor, fill=green!8, draw=green!35, right=3.8cm of mcp] (neo4j) {Neo4j + GDS};

        % === TIMELINES ===
        \draw[tl] (agent.south) -- ++(0,-8.0);
        \draw[tl] (mcp.south) -- ++(0,-8.0);
        \draw[tl] (neo4j.south) -- ++(0,-8.0);

        % === MESSAGES ===
        % 1: Agent -> MCP: tool call
        \node[stepnum, draw=blue!60, fill=blue!8, left=0.15cm] at ([yshift=-1.3cm]agent.south) {1};
        \draw[msgarr, color=blue!60] ([yshift=-1.3cm]agent.south) --
            node[msg, above] {\texttt{get\_relevant\_context(task\_desc)}}
            ([yshift=-1.3cm]mcp.south);

        % 2: MCP -> Neo4j: seed selection
        \node[stepnum, draw=teal, fill=teal!8, left=0.15cm] at ([yshift=-2.5cm]mcp.south) {2};
        \draw[msgarr, color=teal] ([yshift=-2.5cm]mcp.south) --
            node[msg, above] {Entity matching + BM25 query}
            ([yshift=-2.5cm]neo4j.south);

        % 3: Neo4j -> MCP: seeds
        \draw[msgarrdash, color=teal] ([yshift=-3.4cm]neo4j.south) --
            node[msg, above] {Matched seed node IDs}
            ([yshift=-3.4cm]mcp.south);

        % 4: MCP -> Neo4j: PPR
        \node[stepnum, draw=orange!70, fill=orange!8, left=0.15cm] at ([yshift=-4.3cm]mcp.south) {3};
        \draw[msgarr, color=orange!70] ([yshift=-4.3cm]mcp.south) --
            node[msg, above] {\texttt{gds.pageRank.stream(seeds)}}
            ([yshift=-4.3cm]neo4j.south);

        % 5: Neo4j -> MCP: PPR results
        \draw[msgarrdash, color=orange!70] ([yshift=-5.2cm]neo4j.south) --
            node[msg, above] {Ranked nodes + PPR scores}
            ([yshift=-5.2cm]mcp.south);

        % 6: MCP -> Agent: results
        \node[stepnum, draw=blue!60, fill=blue!8, left=0.15cm] at ([yshift=-6.2cm]mcp.south) {4};
        \draw[msgarrdash, color=blue!60] ([yshift=-6.2cm]mcp.south) --
            node[msg, above] {Ranked files + source code (JSON)}
            ([yshift=-6.2cm]agent.south);

        % 7: Agent -> MCP: optional follow-up
        \node[stepnum, draw=gray!50, fill=gray!8, left=0.15cm] at ([yshift=-7.3cm]agent.south) {5};
        \draw[msgarr, color=gray!50] ([yshift=-7.3cm]agent.south) --
            node[msg, above] {\texttt{query\_dependencies(entity)} \textcolor{gray}{\scriptsize[optional]}}
            ([yshift=-7.3cm]mcp.south);

    \end{tikzpicture}
    \caption[Agent--MCP server interaction sequence]{Sequence diagram showing the interaction between an AI agent and the CodeGraph MCP server during a typical retrieval session. \textbf{Step~1:} the agent sends a task description via the \texttt{get\_relevant\_context} tool call. \textbf{Steps~2--3:} the server performs seed selection by matching entity names and running BM25 against the graph. \textbf{Steps~3--4:} the server executes PPR on the GDS projection and collects ranked results. \textbf{Step~4:} post-processed results with source code are returned to the agent. \textbf{Step~5 (optional):} the agent explores dependencies of suspicious entities before generating a solution. Solid arrows indicate requests; dashed arrows indicate responses.}
    \label{fig:agent-workflow}
\end{figure}

This workflow integrates naturally into existing agent frameworks that support MCP, including Claude Code, Cursor, and VS~Code Copilot. The agent requires no knowledge of graph algorithms or code parsing---it simply calls a tool with a task description and receives relevant source code in return.

\subsection{System Prompt and Agent Guidance}
\label{sec:system-prompt}

The MCP server includes a system prompt (defined in \texttt{prompts.py}) that guides agent behavior through standard operating procedures. The prompt instructs the agent to call \texttt{get\_relevant\_context} at the start of each task, to read the returned files before making changes, and to use \texttt{query\_dependencies} when the initial context is insufficient. This prompt is optional---agents can invoke the tools without it---but it improves effectiveness by establishing a consistent retrieval-first workflow.


% =============================================================================
% 3.8 Design Limitations
% =============================================================================
\section{Design Limitations}
\label{sec:design-limitations}

Every design involves trade-offs. The decisions described in this chapter---static-only analysis, Python-specific parsing, conservative call resolution, fixed batch rebuilds---each sacrifice something in exchange for simplicity, correctness guarantees, or implementation tractability. This section makes those trade-offs explicit and quantifies their impact where possible, so that the evaluation results in \cref{ch:evaluation} can be interpreted in the right context.

\paragraph{Static analysis only.} CodeGraph builds its graph from static source code analysis. Dynamic behaviors---monkey-patching, \texttt{eval()}-based code generation, runtime module loading via \texttt{importlib}---are invisible to the parser and produce no graph edges. In codebases that rely heavily on dynamic patterns, such as plugin frameworks or heavily metaclass-driven code, the graph underrepresents the true dependency structure. This limitation is mitigated in practice because the SWE-bench benchmark consists of well-structured, statically analyzable repositories, but it would affect deployment on codebases that use Python's dynamic features extensively.

\paragraph{Python-only parsing.} The current parser implementation supports only Python. The graph schema, PPR algorithm, and MCP server are all language-agnostic, so adding a new language requires only implementing a language-specific parser with entity extraction and call resolution logic---but that is a non-trivial engineering investment. Python-only support means the system cannot be applied to polyglot repositories without modification. SWE-bench is Python-only, so this limitation does not affect the evaluation in \cref{ch:evaluation}, but it would need to be addressed before broader deployment.

\paragraph{Incomplete call resolution.} Without type inference, duck-typed method calls (e.g., \texttt{obj.save()} where the type of \texttt{obj} is unknown) cannot be resolved. The resolver drops ambiguous calls to avoid false edges, which means that between 20\% and 40\% of call sites across typical Python repositories produce no \texttt{CALLS} edge---a range consistent with the prevalence of dynamically typed method dispatch relative to directly resolvable function calls. PPR is robust to missing edges (it propagates through whatever structure is available), but missing edges reduce the reachable set and may exclude relevant files that would otherwise be reached via call chains. This limitation correlates with lower retrieval performance on tasks involving deeply encapsulated code, as discussed in \cref{ch:evaluation}.

\paragraph{Utility function noise.} IDF reweighting (\cref{sec:idf-reweighting}) reduces the influence of highly connected nodes but does not eliminate it entirely. Ubiquitous utility functions can still attract PPR probability and displace more specific results, particularly when seed nodes are in their immediate neighborhood. This effect is visible in the evaluation as a secondary failure mode distinct from the seed quality problem.

\paragraph{No incremental graph updates.} The graph must be fully rebuilt when source code changes, taking 5--30\,seconds depending on repository size. For small repositories (under 10,000 lines), this completes in under 5\,seconds and is barely noticeable. For the largest SWE-bench repositories---which can exceed 100,000 lines across hundreds of files---rebuild latency can reach 30\,seconds, making real-time re-indexing impractical during active development. The two-pass construction and \texttt{MERGE}-based idempotent writes lay the groundwork for future incremental update support, but selective re-parsing of changed files is not yet implemented.


% =============================================================================
% 3.9 Chapter Summary
% =============================================================================
\section{Chapter Summary}
\label{sec:design-summary}

This chapter presented the design of the CodeGraph retrieval pipeline. The system transforms a Python repository into a weighted property graph via tree-sitter parsing and two-pass graph construction, translates a natural language task description into a personalization vector via multi-signal seed selection, ranks all graph nodes by structural relevance via Personalized PageRank, and delivers the top-$k$ results to an AI agent through an MCP server interface.

Three organizing principles guided the design decisions described in this chapter. \emph{Parsimony}: CodeGraph performs only static analysis, avoiding the complexity and runtime cost of dynamic analysis, type inference, or execution tracing---accepting known blind spots (dynamic imports, duck-typed calls) in exchange for simplicity and speed. \emph{Composability}: each stage has a well-defined input and output type, so individual components can be improved, replaced, or extended without affecting the rest of the pipeline, as demonstrated by the iterative seed quality improvements in \cref{ch:evaluation}. \emph{Empiricism}: parameter choices---damping factor, top-$k$, signal weights, token budget---were selected based on measured benchmark performance rather than theoretical defaults. Together, these principles produce a system that is easy to reason about, straightforward to extend, and grounded in observed rather than assumed behavior.

Three specific decisions proved particularly consequential for retrieval quality. First, \emph{uniform PPR restart} outperforms weighted restart because it prevents overconfident seed weights from dominating the computation---a finding that suggests seed diversity matters more than seed confidence when seed quality is imperfect. Second, \emph{IDF-based edge reweighting} reduces the influence of ubiquitous call targets that would otherwise dilute PPR scores across the graph. Third, \emph{compound tokenization} in BM25 seed selection addresses vocabulary fragmentation, a first-order problem in code retrieval. \Cref{ch:evaluation} quantifies the individual and combined impact of each decision, isolating their contributions through a series of controlled ablations.

\Cref{ch:implementation} transitions from design to implementation, presenting the technology stack, code organization, data models, and evaluation harness that transform these architectural principles into a functioning system.